\chapter{Ընդլայնումներ (ընտրովի)}

Ստորև բերվածները դիտվում են որպես \textbf{առաջադեմ ընդլայնումներ}, ոչ պարտադիր MVP պահանջներ.

\begin{itemize}[leftmargin=1.4em]
  \item \textbf{AXI DMA / DDR ուղի} — բարձր տվյալահոսքի դեպքում AXI4-Stream
        $\rightarrow$ AXI DMA $\rightarrow$ DDR $\rightarrow$ PS։
  \item \textbf{FPGA reconfiguration} — Web հարթակից նախապես գեներացված
        կոնֆիգուրացիաների (bitstream-ների) միջև ընտրություն (օր.՝ Environmental /
        Vibration / High-Speed acquisition)։
  \item \textbf{Ֆիզիկական սենսորներ} — էմուլացված աղբյուրների փոխարեն կամ դրանց կողքին։
  \item \textbf{Ադապտիվ հետադարձ կապ (closed loop)} — սերվերի մոդելը հայտնաբերում է
        վարքի փոփոխություն և ավտոմատ փոխում FPGA-ի պարամետրերը (threshold, filter
        window, sample rate, մշակման ռեժիմ)՝ ստեղծելով ադապտիվ ապարատա-ծրագրային համակարգ։
  \item \textbf{Linux Zynq PS-ի վրա} — սերվերային բաղադրիչների մասնակի տեղափոխում Zynq PS-ի Linux-ի վրա։
\end{itemize}

\summarybox{Առաջադրանքը սահմանում է միջամբիոնային դիպլոմային նախագիծ, որտեղ առաջին
ուսանողը կենտրոնանում է FPGA/RTL ապարատի, իսկ երկրորդը՝ ծրագրային հարթակի ու
տվյալների վերլուծության վրա, իսկ ընդհանուր համակարգը ապահովում է երկկողմանի
հաղորդակցություն և runtime կոնֆիգուրացիա։}
